\title{LongRoPE: Extending LLM Context Window Beyond 2 Million Tokens}

\begin{abstract}
The ability to process large context windows is an important attribute for large language models (LLMs). Nonetheless, current methods for expanding context windows remain capped at approximately 128k tokens, constrained by the substantial expenses of fine-tuning, limited availability of extensive textual data, and the disruptive effects caused by introducing novel token positions.
\end{abstract}

This work presents LongRoPE, a method that, for the first time, extends pre-trained LLMs' context window to an unprecedented \textbf{2048k} tokens. Remarkably, this expansion is accomplished with a maximum of only 1k fine-tuning steps at training lengths of up to 256k, all while preserving the model’s original short-context performance. This capability is enabled by three core contributions: (i) the identification and utilization of two types of non-uniformities in positional interpolation via an efficient search process, which facilitates improved fine-tuning initialization and supports up to an 8$\times$ extension even without fine-tuning; (ii) a staged extension approach where the model is initially fine-tuned to a 256k context window, followed by an additional round of positional interpolation applied to this extended LLM, resulting in a 2048k context window; (iii) a recalibration of LongRoPE at 8k context length, ensuring recovery of short context window performance. Extensive empirical evaluations with LLaMA2 and Mistral on a range of tasks validate the approach’s effectiveness. Importantly, models adapted with LongRoPE retain their original architectures with only minor adjustments to the positional embeddings, and the majority of prior optimizations remain compatible. The implementation will be released at \texttt{https://github.com/microsoft/LongRoPE}
\end{abstract}

\section{Introduction}

Large Language Models (LLMs) have achieved impressive results across multiple benchmarks~\cite{gpt4,llama2}, but they are constrained by relatively small context window sizes; for instance, LLaMA2 is limited to 4096 tokens~\cite{llama2}. When sequences exceed this limit, the performance of LLMs drops, as the models have not encountered such positions during training. This limitation creates difficulties in crucial applications such as in-context learning involving many examples~\cite{Huang2023FewerIM} as well as in constructing LLM agents~\cite{park2023generative,madaan2023selfrefine}.

Recent studies have demonstrated that the context window of a pre-trained LLM can be expanded to approximately 128k tokens by leveraging fine-tuning techniques on lengthier input sequences~\cite{longlora,pi,yarn,zhang2024soaring,liu2023scaling}. Nevertheless, further increasing the context window encounters three significant challenges. The first obstacle involves the introduction of untrained position indices, which often leads to an abundance of outlier values, thus causing out-of-distribution complications and impeding the convergence of fine-tuning processes~\cite{pi}. This problem becomes especially severe when extending from 4k to over 1000k tokens, as such a transition introduces more than 90\% unseen positions. The second challenge stems from the reliance of fine-tuning on the availability of suitably lengthy training data. Existing datasets contain relatively few samples with text lengths exceeding 1000k, and training with these ultra-long examples imposes steep computational demands, necessitating significant amounts of GPU time and memory. The third difficulty arises when attempting to stretch the context window to extremely large sizes; under these conditions, attention mechanisms distribute focus across a vast array of token positions, resulting in diluted attention and consequently diminishing the model’s effectiveness on shorter input sequences~\cite{pi}.

A commonly used strategy to address the first challenge involves interpolating RoPE positional embeddings~\cite{rope,pi}, whereby new positional indices are rescaled into the pretraining interval, as depicted in Fig.\ref{fig:overview}. Position Interpolation (PI)~\cite{pi} performs linear interpolation of RoPE's rotational angles according to the ratio of sequence extension. In contrast, NTK~\cite{ntk,dynamicntk} proposes an approach involving non-uniform interpolation and extrapolation schemes across various RoPE dimensions. Meanwhile, YaRN~\cite{yarn} classifies RoPE dimensions based on frequency into three separate categories and then individually applies extrapolation, NTK, or linear interpolation to each group. Nevertheless, positional embeddings in Transformer models possess \emph{complex, uneven distributions of information entropy}. Current methods do not fully exploit this intricate non-uniformity, which results in a loss of informative content and, consequently, constrains the maximum effective context window.

Section~\ref{sec:analysis} demonstrates two main empirical observations: \textbf{(1)} For positional interpolation to be effective, it is essential to address two sources of non-uniformity: differences across RoPE dimensions and across token positions. Lower RoPE dimensions and early token positions require reduced amounts of interpolation, but the best configuration is contingent upon the specific length being extended to. \textbf{(2)} Incorporating these identified non-uniformities into the interpolation strategy enables preservation of vital information present in the original RoPE, especially within critical dimensions and token locations. This approach reduces the degradation induced by interpolation and, as a result, yields a more advantageous starting point for subsequent fine-tuning. Additionally, it permits up to an 8$\times$ extension of context length even without fine-tuning.

Inspired by these insights, we present \textbf{LongRoPE}, a robust technique that enables large language models (LLMs) to operate with context windows surpassing 2 \emph{million} tokens. LongRoPE incorporates three central advancements. Firstly, it capitalizes on multidimensional positional interpolation by accommodating non-uniformities across RoPE dimensions. For every RoPE dimension, LongRoPE determines optimal rescale factors for the rotation angles by taking token positions into account. Given that the space for identifying suitable rescale factors grows exponentially as the extension ratio increases, LongRoPE adopts an evolutionary search framework augmented with two specialized optimization strategies to enhance search performance. An illustration of the rescaled RoPE obtained through this method is provided in Fig.~\ref{fig:overview}.

Subsequently, {LongRoPE} employs a resource-efficient, incremental extension approach to reach a context window of 2048k tokens, circumventing the challenge of fine-tuning directly on extremely long texts, which are both rare and difficult to obtain. This approach initiates by targeting a 256k context length on the base pre-trained LLM, followed by fine-tuning at this specified length. Due to our use of non-uniform positional interpolation, which makes possible an eightfold expansion without requiring further fine-tuning, we next perform an additional search to identify new RoPE rescaling factors on the previously fine-tuned, length-extended LLM. Through this sequence, the 2048k context window is ultimately realized for LLaMA2 and Mistral~\cite{mistral}.

In order to prevent performance loss on the initial (shorter) context window, {LongRoPE} further modifies the RoPE rescale factors within the extended LLM. Employing the same approach used when enlarging from 256k to 2048k, we now downscale to 4k and 8k context lengths on the 256k fine-tuned LLM by leveraging our search algorithm to reduce positional interpolation. During inference, when the input sequence contains fewer than 8k tokens, RoPE is adjusted using the corresponding rescale factors identified through our search.

Through comprehensive evaluations involving multiple LLMs and a range of long-context tasks, our approach has proven to be highly effective. The results indicate that LongRoPE consistently preserves low perplexity across evaluation lengths from 4k up to 2048k, exceeds 90\% passkey retrieval accuracy, and matches the performance of standard benchmarks constructed with a 4096-token context window. Notably, LongRoPE is compatible with all LLMs that utilize RoPE embeddings. Our codebase and the LongRoPE-2048k models will be made publicly available.

\section{Non-uniformity in Positional Interpolation}
\label{sec:analysis}

\subsection{Preliminary}

Transformer architectures inherently lack a mechanism to capture token order and thus depend on additional positional encodings, typically realized as position embeddings. In this study, we investigate the RoPE~\cite{rope} positional embedding method, which is a prevalent choice among modern LLMs. The RoPE encoding for a token located at position $n$ can be expressed in a simplified form as:

\begin{equation}
		\fontsize{7.8}{7.8} \selectfont
	\label{eq:rope}
	\left[ cos(n\theta_0), sin(n\theta_0), cos(n\theta_1), \cdot\cdot\cdot, cos(n\theta_{d/2-1}), sin(n\theta_{d/2-1}) \right]   
\end{equation}

Here, $d$ denotes the dimension of the embedding, $n\theta_i$ corresponds to the rotary position angle for the token at index $n$, and $\theta_i=\theta^{-2i/d}$ defines the set of rotation frequencies. Typically, RoPE uses $10000$ as the default value for the base $\theta$.

\textbf{\emph{Context window extension ratio $s$} and positional interpolation}. We define $s$ as  the ratio of extended context length $L'$ to the original length $L$: $s=\frac{L'}{L}$. 

In order to increase the context window from $L$ to $L'$, existing approaches for positional interpolation propose adjusting the rotation frequencies $\theta_i$ according to the scaling ratio $s$. Defining $\beta = \theta^{2/d}$ and using $\lambda$ as the effective scaling parameter tied to $s$, we can represent all these positional interpolation techniques in a consolidated formulation:

\begin{equation}	
	\fontsize{7.5}{7.5} \selectfont
	\label{eq:unified_rope}
	\begin{aligned}
		\left[cos\left(\frac{n}{\lambda(\beta)^0}\right), sin \left(\frac{n}{\lambda(\beta)^0}\right), cos\left(\frac{n}{\lambda(\beta)^1}\right), \cdots, sin \left(\frac{n}{\lambda(\beta)^{d/2-1}}\right) \right]   
	\end{aligned}
\end{equation}

\textbf{Linear positional interpolation (PI)}. PI~\cite{pi} employs a strategy of linearly interpolating position indices up to the maximum sequence length used during pre-training. Given a desired extension ratio $s$, it uniformly scales down the rotation angles at each position by a factor $\lambda = s$ in every RoPE dimension. Nevertheless, this causes positional encodings to become densely packed, which diminishes the model's capacity to differentiate tokens that are nearby in the sequence. Consequently, PI's effectiveness declines notably when dealing with large extension ratios.

\textbf{NTK-based interpolation and extrapolation}. ~\cite{ntk,dynamicntk} analyze RoPE through the framework of information representation and leverage the Neural Tangent Kernel (NTK) theory~\cite{jacot2018neural,tancik2020fourier}. To address the problem of crowded positions inherent in PI, they propose distributing the effects of interpolation more evenly across the different RoPE dimensions. Specifically, their approach involves reducing the scaling of dimensions corresponding to higher frequencies and increasing it for those associated with lower frequencies. This adjustment enables both positional interpolation and extrapolation, characterized by $\lambda=s^{i}$. The enhanced dynamic NTK method~\cite{dynamicntk} further modulates the extension ratio dynamically at each position, depending on the length of the current input sequence. While PI approaches generally require model fine-tuning, NTK-aware techniques are capable of expanding context lengths without additional training, although practical implementations are often limited to an extension ratio of up to 4$\times$.

\textbf{YaRN}~\cite{yarn} presents a notable enhancement in positional interpolation capabilities. The method segments the RoPE dimensions into three separate clusters according to frequency, applying distinct interpolation methods for each. Specifically, dimensions with higher frequencies are subjected to extrapolation ($\lambda$=1), while those associated with lower frequencies utilize linear interpolation (PI). Dimensions that fall between these two groups make use of the NTK approach. A crucial feature of YaRN is its frequency-based grouping of RoPE dimensions, a process which, at present, relies on empirically driven human judgment. As a consequence, this manual grouping may limit performance when applied to unfamiliar LLMs.

\subsection{Study on Non-uniform Positional Interpolation}

Drawing on insights from NTK and YaRN, we observe that their improvements stem from leveraging non-linear mechanisms, particularly through accounting for frequency variations along RoPE dimensions to enhance interpolation and extrapolation capabilities. Nevertheless, most non-linear approaches in use today depend significantly on heuristics crafted by humans. This prompts two central inquiries: (1) Does the existing strategy for positional interpolation achieve optimality? (2) Could there exist additional non-linearities that have yet to be investigated?

In addressing these questions, we employ evolution search (refer to Sec.~\ref{sec:method}) to identify more effective non-uniform positional interpolations for LLaMA2-7B. The process is steered by perplexity metrics, evaluated on five randomly selected passages from the PG19~\cite{pg19} validation set. Our empirical investigation uncovers several important insights.

\textbf{Finding 1}: \textit{RoPE dimensions exhibit substantial non-uniformities, which are not effectively handled by current positional interpolation methods.}

We optimize $\lambda$ for each RoPE dimension in Eq.~\ref{eq:unified_rope}. Table~\ref{tbl:exp1} reports the perplexity of LLaMA2-7B using various approaches on the PG19 and Proof-pile~\cite{proof-pile} test sets, evaluated without additional fine-tuning. Our method achieves substantial perplexity reductions, indicating that prevalent linear (PI) and existing non-uniform methods (Dynamic-NTK and YaRN) do not reach optimal performance. Interestingly, YaRN performs worse than both PI and NTK on the PG19 benchmark, as its extrapolation does not sufficiently extend the usable context window in models that have not been further fine-tuned; as a consequence, perplexity for YaRN increases abruptly after 7k tokens in an 8k context window.

In our investigation, the rescaling coefficients $\lambda$ in Eq.~\ref{eq:unified_rope} are non-uniform, in contrast to the constant scaling factor $s$ applied in PI, the expression used in NTK, and YaRN's group-based approach. Utilizing these varying $\lambda$ values substantially enhances LLaMA2's language modeling ability (perplexity) for context lengths of 8k and 16k, all without additional fine-tuning. The reason behind this improvement lies in how the modified positional embeddings better retain the structure of the original RoPE, particularly within crucial dimensions, thereby making it easier for the LLM to distinguish among tokens that are positioned closely together.

\textbf{Finding 2}: \textit{RoPE for the initial tokens in the input sequence should be extrapolated with less interpolation.} 

We propose that position interpolation for the initial $\hat{n}$ tokens in input sequences should be minimized, as these tokens typically obtain higher attention scores and thus play a pivotal role within attention layers. This hypothesis aligns with findings from Streaming LLM~\cite{streamingllm} and LM-Infinite~\cite{han2023lminfinite}. To test this, we expand the context window to 8k and 16k tokens by applying PI and NTK, respectively, while excluding the first $\hat n$ tokens ($\hat n$ taking values of 0, 2, ..., 256) from interpolation. Notably, when $\hat n$ is set to 0, the method defaults to standard PI and NTK implementations. Table~\ref{tbl:tokenposition} reveals two principal results: \textbf{(1)} excluding initial tokens from position interpolation consistently yields performance gains; and \textbf{(2)} the choice of $\hat n$ that provides optimal results varies with the length of the context window being extended.

\textbf{Finding 3}: \textit{Non-uniform positional interpolation effectively extends LLM context window in both fine-tuning and non-fine-tuning settings.} 

Although we have demonstrated that employing our discovered non-uniform position interpolation yields notable gains in extension capabilities at 8k and 16k context lengths without the need for fine-tuning, extending to even longer contexts necessitates a fine-tuning stage. Therefore, we fine-tune LLaMA2-7B with our optimized RoPE configuration for a 64k context window (refer to the Appendix for details). According to Table~\ref{tbl:finetune}, our approach significantly surpasses both PI and YaRN, not only prior to fine-tuning but also afterward. This performance edge arises from our strategic use of non-uniform positional interpolation, which reduces the risk of information degradation and offers an advantageous initialization point for subsequent fine-tuning.

\textbf{Summary}. Our work identifies two sources of non-uniformity: heterogeneity across RoPE dimensions and variation among token positions. Effectively leveraging these non-uniformities within positional interpolation results in considerable enhancement of context extension performance for LLMs.

\section{LongRoPE}

\label{sec:method}
Building on these observations, we propose LongRoPE, a method that initially incorporates a fast search algorithm to leverage both forms of non-uniformity. This approach subsequently enables LLMs to process context windows exceeding 2 million tokens.

\subsection{Problem Formulation}

The presence of these two types of non-uniformity results in an expansive solution space and increases the difficulty of optimization. To manage this, we recast the multidimensional non-uniform position interpolation task as a search problem.

Consider an LLM configured for a context window of size $L'$ and tasked with processing extended input documents $\mathbf{X}$, with each segment $\mathbf{x}\in\mathbf{X}$ exceeding $L'$ tokens in length. We represent the initial rotary angle for the $i^{th}$ dimension in the RoPE embedding at the $n^{th}$ token position as $\frac{n}{\beta^i}$. The optimization objective is expressed as:

\begin{equation}
\begin{gathered}
		\label{eq:problem}

\underset {\mathbf{x}\in\mathbf{X}; \, |\mathbf{x}|\ge L'}{ \operatorname {arg\,min} } \,  \mathcal{L}\, \left( \text{LLM} (\text{RoPE}, \mathbf{X})\right),	\text{where \;} 
\\
\scalebox{0.93}{
$\underset {\begin{subarray}{l} i=0,\cdot\cdot,\frac{d}{2}-1;\\ n\in[ 0,|\mathbf{x}|);\end{subarray}}{\text{RoPE}_(n)}= {\left[\cdots, \cos\Bigl(\mathbb{I}(\hat{\lambda}_i, \hat{n})\frac{n}{\beta^i}\Bigr), \sin\Bigl(\mathbb{I}(\hat{\lambda}_i, \hat{n})\frac{n}{\beta^i}\Bigr), \cdots\right]}$}\\
	\text{where \;} \mathbb{I}(\hat{\lambda}_{i},\hat{n})=
\begin{cases}
	1&\text{if $n<\hat{n}$}\\
	{\frac{1}{\lambda}_{i}}&\text{if $n\geq\hat{n}$}
\end{cases}
	\end{gathered}
\end{equation}

Here, we define a collection of scaling factors, $\mathbb{I}(\hat{\lambda}_{i},\hat{n})$, designed to accommodate both categories of non-uniformities. The parameters $\hat{\lambda_i}$ and $\hat{n}$ signify, respectively, the degree of non-uniformity in RoPE dimensions and token position indices. In practice, $\mathbb{I}(\hat{\lambda}_{i},\hat{n})$ functions to adjust the rotation angle specific to the $i^{th}$ RoPE dimension: $\hat\lambda_{i}$ serves as the scaling factor, and $\hat{n}$ sets the threshold for token positions. For the first $\hat{n}$ tokens ($n<\hat{n}$), the scaling factor $\hat\lambda_{i}$ remains inactive, and the standard RoPE rotation angle $\frac{n}{\beta^i}$ is maintained. Once token positions reach $n\ge\hat{n}$, the scaling factor is implemented.

With a desired context window length of $L'$, the task is to determine the most effective set of rescale factors ($\mathbb{I}(\hat{\lambda}_{0},\hat{n})$, $\mathbb{I}(\hat{\lambda}_{1},\hat{n})$, ...$\mathbb{I}(\hat{\lambda}_{i},\hat{n})$...) corresponding to each RoPE dimension from the $1^{st}$ through the $d^{th}$. The aim is for the target $\text{LLM}$—utilizing the rescaled $\text{RoPE}$—to achieve the lowest possible next token prediction loss, $\mathcal{L}$ (also known as perplexity), for input sequences $\mathbf{X}$ with token length $L'$.

\subsection{Searching the Non-uniform Position Interpolation}

To address the challenge presented in Eq.~\ref{eq:problem}, we propose an intuitive yet powerful approach that systematically explores the optimal set of RoPE rescale factors, thereby capitalizing on the multidimensional non-uniformities inherent to position embeddings.

\textbf{Search space}. Our method constructs a comprehensive search space designed to incorporate both sources of non-uniformity, as summarized in Table~\ref{tbl:searchspace}. In particular, we facilitate the assignment of distinct rescale factors for each individual dimension of the RoPE embedding. To simplify the structure of the search space, the algorithm searches over $\lambda_i$ and $\hat{n}$, as opposed to directly optimizing $\mathbb{I}(\hat{\lambda}_{i},\hat{n})$, using the relationship $\hat{\lambda}_{i}=1/\lambda_i$. Referring to Table~\ref{tbl:searchspace}, $\lambda_i$ is chosen within a continuous range: the minimum possible value is 1.0 (corresponding to extrapolation), and the maximum is $s\times1.25$ (enabling broader interpolation than Position Interpolation, PI), using increments of 0.01. Here, $s$ denotes the intended context window extension ratio.

The parameter $\hat{n}$ determines how many of the starting token positions are preserved with the original RoPE embedding, foregoing positional interpolation. In practice, $\hat{n}$ is selected from the set \{0, 1, 2, 4, 8, 12, 16, 20, 24, 28, 32, 64, 128, 256\}. Notably, if $\hat{n}=0$, then rescale factors found by search are applied to every token position.

\textbf{Evolution-based search}. The search space detailed in Table~\ref{tbl:searchspace} encompasses a vast array of possible positional interpolation configurations, rendering exhaustive exploration highly impractical. For instance, achieving a $s=4\times$ extension yields $400^{128/2}\times14 = 4\times10^{167}$ possible options. As the extension ratio grows, the number of potential configurations increases exponentially. To manage this complexity, we adopt an evolution search strategy~\cite{guo2020single} and employ two optimization methods that significantly enhance the efficiency of the search. The entire procedure is outlined in Algorithm~\ref{alg:search}.

\textit{Optimized initial population generation}. Rather than solely relying on random assignment for the initial population of $P$ rescale factors, we explicitly incorporate the three RoPE rescale factors associated with PI, NTK, and YaRN as members of this population from the outset. For the other $P$-3 individuals, each of the three rescale factors undergoes mutation at a probability of $p$, resulting in their random variation.

\textit{Monotonically non-decreasing constraint}. Following the creation of the initial population, we evaluate the perplexity of the LLM for each individual. To do this, the relevant RoPE rescale factors are applied to the target LLM, after which the perplexity is determined using the input $\mathbf{X}$. The individuals with the top-k scores are then chosen as the parent set for subsequent evolutionary steps. Nonetheless, given the high dimensionality of the search space, naive application of mutation and crossover often results in the examination of suboptimal candidates, thus incurring excessive perplexity evaluations. This issue is exacerbated when $L'$ is large, since each perplexity inference requires a considerable amount of computational time.

To tackle this issue, we enforce a monotonic non-decreasing constraint on the selected RoPE rescaling factors so that $\lambda_i\le \lambda_{i+1}$. We only utilize RoPE configurations adhering to this monotonicity when conducting perplexity measurements with LLMs, thereby greatly limiting the computational expense of the search. More concretely, we stipulate that $\lambda_i$ must monotonically increase as the RoPE dimension $i$ progresses (for $i=0,\ldots,63$). This requirement is motivated by NTK theory~\cite{jacot2018neural,tancik2020fourier,ntk}, which posits that lower-dimensional components (characterized by higher frequencies) benefit from minimal interpolation (smaller $\lambda_i$ values), whereas higher-dimensional components (with lower frequencies) can accommodate greater degrees of interpolation (larger $\lambda_i$ values).

\begin{algorithm}[t]
	\small
	\caption{The search algorithm}
	\textbf{Input:} target LLM, input samples $\mathbf{X}$,   population size $P$, mutation size $N_1$, crossover size $N_2$, max iterations $\mathcal{T}$, mutate probability $p$\\

	\begin{algorithmic}[1]
		\label{alg:search}
		\STATE Top-k $\gets \phi$;	
		\STATE $\text{P}_0 \gets$ \textit{Initialize\_population\_with\_optimization} ($P$, $\mathbf{X}$, $p$); 
		\FOR{$i = 1$ to $\mathcal{T}$}
		\STATE \textit{Compute\_perplexity} (LLM, $\text{P}_{i-1}$, $\mathbf{X}$);
		\STATE  Top-k $\gets$ \textit{Update\_Topk} (Top-k, $\text{P}_{i-1}$);
		\STATE$\text{P}_{mutation} \gets$ \textit{Mutation\_with\_mono\_constraint} (Top-k, $N_1$, $p$); 
		\STATE $\text{P}_{crossover} \gets$ \textit{Crossover\_with\_mono\_constraint} (Top-k, $N_2$);
		\STATE $\text{P}_i = \text{P}_{mutation} \cup \text{P}_{crossover} \cup \text{Top-k}$;
		\ENDFOR
		\STATE Return the member of Top-k with the minimum perplexity;
	\end{algorithmic}
\end{algorithm}

\textbf{8$\times$ extension without fine-tuning}. By leveraging evolutionary search, our approach discovers optimal non-uniform RoPE rescale factors that retain crucial dimensional and positional information, thereby reducing information loss typically caused by interpolation. As shown in Fig.\ref{fig:non-ft}, this strategy allows us to expand LLaMA2's context window from 4k to 32k tokens without requiring any fine-tuning. By comparison, current techniques like PI, along with non-uniform NTK and YaRN, exhibit a sharp increase in perplexity when the window is extended beyond 2$\times$.

\subsection{Extending LLM Context Window to 2048K}
\label{sec:extendto2048k}

\textbf{Progressive extension to 2048k}. In this section, we present our strategy for expanding the context window of pre-trained LLMs from the conventional 4k length up to beyond 2048k. Our results indicate that the use of non-uniform positional interpolation permits an extension by a factor of 8 without the need for fine-tuning. However, achieving more substantial increases (such as a 512$\times$ enlargement) necessitates fine-tuning. A common approach involves identifying suitable RoPE rescaling factors that correspond to the target length of 2048k, followed by a fine-tuning process. Nevertheless, this strategy is hindered by the substantial computational demands involved. Additionally, our observations suggest that effective fine-tuning of LLMs at such large extension scales remains a significant challenge (refer to Appendix).

Luckily, LongRoPE demonstrates strong performance with both the base and further fine-tuned extended LLMs. Accordingly, we propose an efficient and incremental approach that attains the desired 2048k target using merely 1k fine-tuning iterations, all within a training sequence length capped at 256k.

$\Diamond$ \textit{Exploring LongRoPE-based extension of pre-trained LLMs to 256k}. Using LLaMA2 as a case study, we investigate the search process for expanding the context window to 128k and 256k tokens. In this process, the extension ratios correspond to 32$\times$ and 64$\times$, respectively.

$\Diamond$ \textit{Fine-tuning to 256k}. Subsequently, we adapt the pre-trained LLM to support a context window length of 256k through a staged fine-tuning process. Initially, we fine-tune LLaMA2 for 400 iterations with the RoPE rescaling factors set for 128k. Following this, we switch the RoPE rescaling factors to accommodate 256k at the resulting checkpoint and continue with an extra 600 fine-tuning steps. This incremental strategy has demonstrated greater efficiency compared to directly fine-tuning for a 256k context window from the outset.

$\Diamond$ \textit{Expanding a fine-tuned extended LLM to 2048k via LongRoPE search}. In the concluding step, we conduct an additional search on the previously fine-tuned LLM set at 256k in context length. This process allows us to achieve a substantial context window of 2048k, all without necessitating further fine-tuning. As a result, the achieved extension ratio is $512\times$.

\textbf{Performance degradation within shorter context windows}. Upon expanding the context window to an extensive length of 2048k, we observe diminished performance in the initial context span. This phenomenon is well-documented in the context of positional interpolation~\cite{pi}, which constrains the position embeddings in the lower-dimensional subspace corresponding to the original window, leading to a more compressed distribution and, consequently, suboptimal language model output. Specifically, using a 512$\times$ extension ratio causes positions inside the initial 4k window to become especially dense.

To address this, we conduct an additional evolutionary search on the expanded LLM to fine-tune the RoPE rescale factors for shorter context lengths, such as 4k and 8k tokens. Since positional interpolation is less critical at these lengths, we limit the maximum value of $\lambda$ explored during the search. At inference time, the LLM automatically modifies the relevant RoPE rescale factors based on the input context.

\section{LongRoPE}

\label{sec:method}
Building upon these insights, we propose LongRoPE, which begins by implementing an effective search strategy designed to take advantage of the identified dual non-uniformities. This approach is subsequently utilized to expand the LLM context window to exceed 2 million tokens.

\subsection{Problem Formulation}

The presence of these two non-uniformities results in an expansive solution space and adds layers of difficulty to the optimization process. To tackle this challenge, we recast the multidimensional non-uniform position interpolation as a search-based optimization task.

For large language models aiming to process a context window of size $L'$ with long input documents $\mathbf{X}$, such that each segment $\mathbf{x} \in \mathbf{X}$ exceeds $L'$ in token count, we represent the initial rotary position angle for the $i^{th}$ dimension in the RoPE embedding at token index $n$ as $\frac{n}{\beta^i}$. Under these conditions, the optimization task is defined as:

\begin{equation}
\begin{gathered}
		\label{eq:problem}

\underset {\mathbf{x}\in\mathbf{X}; \, |\mathbf{x}|\ge L'}{ \operatorname {arg\,min} } \,  \mathcal{L}\, \left( \text{LLM} (\text{RoPE}, \mathbf{X})\right),	\text{where \;} 
\\
\scalebox{0.93}{
$\underset {\begin{subarray}{l} i=0,\cdot\cdot,\frac{d}{2}-1;\\ n\in[ 0,|\mathbf{x}|);\end{subarray}}{\text{RoPE}_(n)}= {\left[\cdot\cdot,\cos\left(\mathbb{I}(\hat{\lambda}_{i}, \hat{n})\times\frac{n}{\beta^i}\right),  \sin\left(\mathbb{I}(\hat{\lambda}_{i}, \hat{n})\times\frac{n}{\beta^i}\right),\cdot\cdot\right] }$}\\
	\text{where \;} \mathbb{I}(\hat{\lambda}_{i},\hat{n})=\begin{cases}
	1&\mbox{\;  $n<\hat{n}$}\\
{\frac{1}{\lambda}_{i}}&\mbox{\; $n\geq\hat{n}$}
\end{cases}
	\end{gathered}
\end{equation}

In this approach, we incorporate a set of rescale factors, $\mathbb{I}(\hat{\lambda}_{i},\hat{n})$, to address both identified forms of non-uniformity. Here, $\hat{\lambda_i}$ reflects the degree of non-uniformity present in RoPE dimensions, while $\hat{n}$ corresponds to variations across token positions. The function $\mathbb{I}(\hat{\lambda}_{i},\hat{n})$ specifically rescales the rotational angle in the $i^{th}$ dimension of RoPE, employing $\hat\lambda_{i}$ as the scaling parameter and $\hat{n}$ as a position threshold. For the first $\hat{n}$-$1$ tokens, the scaling factor $\hat\lambda_{i}$ is not utilized, and computation proceeds with the standard RoPE rotational angle $\frac{n}{\beta^i}$. Once the position index meets or exceeds $n \ge \hat{n}$, the rescale factor comes into effect.

Suppose the desired context window size is $L'$. The task is to determine the set of optimal rescale factors ($\mathbb{I}(\hat{\lambda}_{0},\hat{n})$, $\mathbb{I}(\hat{\lambda}_{1},\hat{n})$, ..., $\mathbb{I}(\hat{\lambda}_{i},\hat{n})$, ...) for each RoPE dimension from the $1^{st}$ to $d^{th}$. With these rescaled $\text{RoPE}$ parameters, the target $\text{LLM}$ aims to achieve the lowest possible next token prediction loss, $\mathcal{L}$ (measured as perplexity), across input sequences $\mathbf{X}$ of length $L'$.

\subsection{Searching the Non-uniform Position Interpolation}

To address the challenge outlined in Eq.~\ref{eq:problem}, we present a straightforward and effective approach that seeks the optimal RoPE rescale factors in order to capitalize on the multidimensional non-uniformities present in position embeddings.

\textbf{Search space}. Our method employs an expansive search space that accommodates both types of non-uniformity. Table~\ref{tbl:searchspace} summarizes the specifics of this search space. In particular, we permit an independent rescale factor to be discovered for every dimension in the RoPE embedding. To streamline the configuration of the search space, we opt to search over $\lambda_i$ and $\hat{n}$ rather than directly over $\mathbb{I}(\hat{\lambda}_{i},\hat{n})$, where $\hat{\lambda}_{i}=1/\lambda_i$. As detailed in Table~\ref{tbl:searchspace}, $\lambda_i$ is selected from values beginning at 1.0 (corresponding to direct extrapolation) up to $s\times1.25$ (allowing for more extensive interpolation than Position Interpolation), incremented by steps of 0.01, with $s$ denoting the intended ratio for extending the context window.

The parameter $\hat{n}$ determines how many of the earliest token positions utilize the unaltered RoPE embedding, foregoing position interpolation. In practice, $\hat{n}$ is selected from the set \{0, 1, 2, 4, 8, 12, 16, 20, 24, 28, 32, 64, 128, 256\}. If $\hat{n}=0$, then every token position employs the optimized rescale factors.

\textbf{Evolution-based search}. The search space detailed in Table~\ref{tbl:searchspace} encompasses a vast array of positional interpolation strategies, making comprehensive exploration highly complex. For instance, when using a $s=4\times$ extension, there exist $400^{128/2}\times14 = 4\times10^{167}$ possible configurations. As the extension ratio increases, the search space grows at an exponential rate. To effectively navigate this expansive space, we adopt an evolution search approach as described by~\cite{guo2020single}, and implement two optimization strategies designed to substantially improve the efficiency of the search. The main steps of the search algorithm are outlined in Algorithm~\ref{alg:search}.

\textit{Optimized initial population generation}. Rather than populating the initial $P$ rescale factors solely at random, we deliberately include the three RoPE rescale factors associated with PI, NTK, and YaRN as members of the initial population. The other $P$-3 individuals are generated by applying random mutations to these three rescale factors, with each mutation occurring at probability $p$.

\textit{Monotonically non-decreasing constraint}. Following the creation of the initial population, we evaluate the LLM perplexity associated with each candidate. This involves implementing their respective RoPE rescale factors within the target LLM and assessing the perplexity on input $\mathbf{X}$. The highest-ranking $k$ individuals are then selected as parents for subsequent evolutionary steps. Nonetheless, due to the enormous size of the search space, straightforward application of mutation and crossover frequently leads to suboptimal solutions, resulting in an excessive number of perplexity evaluations. This inefficiency is accentuated for large $L'$, as each perplexity computation demands significant inference time.

To tackle this issue, we enforce that the rescaled RoPE factors follow a non-decreasing sequence: $\lambda_i \le \lambda_{i+1}$. This monotonicity constraint is imposed such that only those RoPE variants fulfilling the condition are considered for perplexity assessment in LLMs, thereby greatly curbing the computational expense of the search. Concretely, we stipulate that $\lambda_i$ must strictly increase along the RoPE dimension ($i = 0, \ldots, 63$). This monotonic relationship between dimensions is motivated by findings from NTK theory~\cite{jacot2018neural,tancik2020fourier,ntk}, which indicate that lower-index (high-frequency) components warrant less interpolation (i.e., smaller $\lambda_i$ values), whereas higher-index (low-frequency) components can accommodate increased interpolation (i.e., larger $\lambda_i$ values).

\begin{algorithm}[t]
	\small
	\caption{The search algorithm}
	\textbf{Input:} target LLM, input samples $\mathbf{X}$, population size $P$, mutation size $N_1$, crossover size $N_2$, maximum number of iterations $\mathcal{T}$, mutation probability $p$\\

	\begin{algorithmic}[1]
		\label{alg:search}
		\STATE Top-k $\gets$ $\phi$;
		\STATE $\text{P}_0$ $\gets$ \textit{Initialize\_population\_with\_optimization}($P$, $\mathbf{X}$, $p$);
		\FOR{$i=1$ \textbf{to} $\mathcal{T}$}
		\STATE \textit{Compute\_perplexity}(LLM, $\text{P}_{i-1}$, $\mathbf{X}$);
		\STATE Top-k $\gets$ \textit{Update\_Topk}(Top-k, $\text{P}_{i-1}$);
		\STATE $\text{P}_{mutation}$ $\gets$ \textit{Mutation\_with\_mono\_constraint}(Top-k, $N_1$, $p$);
		\STATE $\text{P}_{crossover}$ $\gets$ \textit{Crossover\_with\_mono\_constraint}(Top-k, $N_2$);
		\STATE $\text{P}_i$ $\gets$ $\text{P}_{mutation}$ $\cup$ $\text{P}_{crossover}$ $\cup$ Top-k;
		\ENDFOR
		\STATE Output the member of Top-k with minimum perplexity;
	\end{algorithmic}
\end{algorithm}

\textbf{8$\times$ extension without fine-tuning}. By employing evolutionary search, our approach successfully discovers optimal non-uniform RoPE rescaling parameters, which retain crucial dimensions and positions, thereby reducing information degradation from interpolation. As shown in Fig.\ref{fig:non-ft}, this enables our technique to expand LLaMA2’s context window from 4k up to 32k tokens without requiring any fine-tuning. Meanwhile, alternative strategies such as PI, as well as non-uniform NTK and YaRN, exhibit a sharp increase in perplexity after merely doubling the window size.

\subsection{Extending LLM Context Window to 2048K}
\label{sec:extendto2048k}

\textbf{Progressive extension to 2048k}. Here, we present our approach for scaling the context window of pre-trained LLMs from the standard 4k up to more than 2048k. Our proposed non-uniform positional interpolation enables an 8$\times$ context extension without the requirement of fine-tuning. However, achieving substantially larger extensions (e.g., 512$\times$) necessitates fine-tuning procedures. A common technique involves searching for RoPE rescaling factors applicable to the targeted 2048k context and then performing fine-tuning. Nonetheless, this strategy encounters significant obstacles due to the immense computational costs associated with training. Furthermore, our practical observations indicate that successfully fine-tuning LLMs for such substantial extension ratios is quite challenging (refer to Appendix).

Fortunately, LongRoPE demonstrates efficacy with both base and further-trained extended LLMs. To this end, we propose an efficient and incremental approach that reaches the desired 2048k context window using only 1k fine-tuning steps at a maximum training length of 256k.

$\Diamond$ \textit{LongRoPE search for extending pre-trained LLMs to 256k.} Using LLaMA2 as a reference model, we perform a search targeting context window lengths of 128k and 256k tokens. At these stages, the corresponding extension factors are 32$\times$ and 64$\times$, respectively.

$\Diamond$ \textit{Fine-tuning to 256k}. Subsequently, we adapt the pre-trained LLM to support a 256k context window via additional fine-tuning. Our procedure involves first fine-tuning LLaMA2 for 400 steps, employing the RoPE rescaling factors set to 128k. Upon completing this phase, we update the RoPE rescaling factors to 256k on the resulting checkpoint and proceed with a further 600 fine-tuning steps. This staged approach demonstrates improved efficiency compared to initiating fine-tuning with the 256k setting from the start.

$\Diamond$ \textit{Scaling fine-tuned extended LLM to 2048k using LongRoPE search}. In the final step, an additional search is conducted on the previously fine-tuned 256k-length LLM. This approach successfully achieves a substantially increased context window of 2048k, eliminating the need for additional fine-tuning. This process yields an extension ratio of $512\times$.

\textbf{Restoration of performance in shorter context windows}. Following the extension of the context window to a considerable 2048k length, there is an observable decline in performance within the model's original context range. This phenomenon is a recognized limitation of positional interpolation~\cite{pi}; by requiring position embeddings for higher-index tokens in the original context span to be confined within a substantially smaller representational space, the resulting compression undermines the model's effectiveness. When utilizing an extension ratio of 512$\times$, the density of positions in the initial 4k context window increases markedly, exacerbating this issue.

To address this, an additional evolution search is conducted on the expanded LLM to fine-tune the RoPE rescale factors for shorter context lengths, such as 4k and 8k tokens. We decrease the upper bound of the searched $\lambda$ values since shorter contexts demand less positional interpolation. At inference time, the LLM automatically modifies the associated RoPE rescale factors as needed.

 \section{Experiments}

\label{sec:exp}
\subsection{Setup}
\textbf{Evaluation Tasks and models}. LongRoPE is implemented with LLaMA2-7B and Mistral-7B, and we assess its effectiveness across three dimensions: (1) perplexity of large language models equipped with extended context lengths when processing lengthy documents; (2) performance on the Passkey retrieval task, which tests a model’s capacity to locate a straightforward passkey amid extensive irrelevant content; and (3) results on standard LLM benchmarks constrained to a context window of 4096 tokens.

\textbf{Fine-tuning}. For LLaMA2, the model is fine-tuned with a learning rate set to 2e-5, implementing a linear decay schedule and a global batch size of 32. Training is conducted for 400 steps on the Redpajama~\cite{redpajama} dataset, which is divided into 128k-length segments marked by BOS and EOS tokens at their boundaries. After completing these steps and saving the checkpoint, the context window is extended to 256k through an additional 600 steps of training. Training on the 128k context size utilizes 8 A100 GPUs, leveraging a distributed framework~\cite{cube}, while scaling up to 256k requires 16 A100 GPUs. For the Mistral model, a fixed learning rate of 1e-6 is used together with a global batch size of 64. The configuration for both the 128k and 256k models follows the protocol established in YaRN~\cite{yarn}, specifically 400 training steps on the Together Computer's Long-Data Collections~\cite{mistraldata} with a 16k sequence length, using 4 A100 GPUs for each run.

\textbf{Search}. For cases where the target window size does not exceed 256k, the configuration parameters are set as follows: $P$=64, $N_1$=$N_2$=16, $p$=0.3, and $\mathcal{T}$=40. During each search iteration, the top-32 candidates are chosen to undergo mutation and crossover. Perplexity is evaluated on five randomly selected samples from the PG19 validation set, ensuring each meets the required minimum target context length. For window sizes greater than 512k, the parameters for population, mutation, and crossover are each reduced by half. In these cases, perplexity assessment uses three randomly chosen samples from the Pile-Books3~\cite{pile} validation set.

\textbf{Baselines}. For expansion to 2048k context windows, we performed fine-tuning on models with either 128k or 256k context lengths, resulting in LongRoPE-2048k (ft=128k) for LLaMA2 and LongRoPE-2048k (ft=256k) for Mistral. We evaluate these four models against leading context extension approaches in the field, focusing on open-source LLMs that have been fine-tuned following positional interpolation methods such as PI, NTK, and YaRN. The comparison set includes models like Together-32k~\cite{together}, Code LLaMA~\cite{codellama}, LongLoRA-full-FT-100k~\cite{longlora}, as well as YaRN-LLaMA and YaRN-Mistral~\cite{yarn}.

\subsection{Main Results}

\label{sec:mainresult}
\textbf{Long sequence language modeling within 256k}. We start by assessing our approach against leading extended LLMs in the context of sequence modeling capped at 256k tokens. To establish the general applicability of our method, we use the Proof-pile~\cite{pg19} and PG19~\cite{pile} datasets, specifically employing their test splits. Perplexity is measured across different context lengths through the use of a 256-token sliding window. In the case of PG19, we utilize the entire set of 100 test documents. For Proof-pile, aligning with the procedure from YaRN~\cite{yarn}, we randomly sample 10 documents, each with a minimum length of 128k tokens.

Table~\ref{tbl:proof-pile} and Table~\ref{tbl:pg19} present a comparison of perplexity scores for LLaMA2 and Mistral, both augmented through various interpolation strategies, on Proof-pile and PG19 datasets, respectively. We emphasize two principal findings: \textbf{(1)} the extended models consistently demonstrate decreasing perplexity as evaluation length increases from 4k to 256k tokens, indicating effective utilization of extended context. \textbf{(2)} Remarkably, even when operating with a context window up to 16$\times$ longer—a scenario that typically poses difficulty for retaining strong performance at shorter lengths—our LongRoPE-2048k models surpass the strongest existing baselines within the 256k context length regime.

\textbf{Long sequence language modeling beyond 2000k}. To assess performance on exceptionally lengthy documents, we utilize the Books3~\cite{pile} dataset. For the sake of efficient evaluation, a random sample of 20 books—each with a length greater than 2048k—is chosen, employing a sliding window of 256k tokens.

According to Table~\ref{tbl:books3}, LongRoPE enables both LLaMA2-7B and Mistral-7B models to reach a context window of 2048k, all while retaining or surpassing baseline perplexity scores at more moderate sequence lengths ranging from 8k to 128k. It is apparent that there are significant variations in performance between the 2048k versions of LLaMA2 and Mistral. Specifically, Mistral demonstrates superiority over baselines at shorter sequence lengths, yet its perplexity climbs above 7 once the context length surpasses 256k. In the case of LLaMA2, the model exhibits the anticipated trend—perplexity steadily decreases as context grows, with only slight increases observed at 1024k and 2048k. Furthermore, LongRoPE-2048k achieves improved results on LLaMA2 when fine-tuned at 256k compared to 128k, attributed to a lower secondary extension ratio (specifically, 8$\times$ rather than 16$\times$). Conversely, Mistral achieves better performance when fine-tuned with a 128k context window. This outcome is primarily due to the fact that, following YaRN's procedure, Mistral is fine-tuned at 128k and 256k using a 16k training length, a factor that limits the model's effectiveness in extending the context window after fine-tuning.

\textbf{Passkey retrieval}. In this section, we investigate the usable context window size for generation tasks. Adopting the synthetic passkey retrieval benchmark introduced by~\cite{passkey}, we evaluate model performance in extracting a randomly selected passkey (a five-digit numeric string) embedded within an extended document. The prompt structure is outlined in the appendix. For evaluation, we conduct 10 runs of the passkey retrieval task, each time positioning the passkey uniformly at random within the given context window.

Fig.~\ref{fig:passkey} illustrates the comparison of retrieval accuracy against baseline approaches. The accuracy of prior models declines sharply to zero beyond the 128k mark. By contrast, even under the stringent challenge of extracting a passkey from a sequence containing millions of tokens, our LongRoPE-LLaMA2-2048k (ft=256k) consistently sustains high retrieval accuracy ($\ge$90\%) throughout the range from 4k up to 2048k tokens. LongRoPE-Mistral-2048k (ft=128k) preserves perfect accuracy up to 1800k, before dropping to 60\% accuracy at 2048k. This decline is consistent with the findings in Table~\ref{tbl:books3}, where a mild increase in perplexity is observed at the 2048k length.

\textbf{Performance on standard benchmarks using the initial context window}. We assess the LongRoPE-2048k models within their original context window by employing the Hugging Face Open LLM Leaderboard~\cite{huggingfaceboard}, considering both zero-shot and few-shot evaluation settings. The specific benchmarks used are ARC-Challenge~\cite{allenai:arc} with 25-shot evaluation, HellaSwag~\cite{zellers2019hellaswag} with 10-shot, MMLU~\cite{mmlu} with 5-shot, and TruthfulQA~\cite{lin2021truthfulqa} in the 0-shot scenario.

As indicated in Table~\ref{tbl:commonsense}, our models perform similarly to the original baseline on the benchmark created for a shorter context window, and actually surpass the original Mistral on TruthfulQA with an improvement of +0.5\%. Although LongRoPE-LLaMA2-2048k—fine-tuned at 256k—exhibits marginally greater performance drops, it still maintains acceptable results across the majority of tasks.

\subsection{Ablation Results}

\textbf{Effectiveness of the second positional interpolation}. Within our progressive extension approach, we implement our search-based algorithm to perform a second round of non-uniform positional interpolation on large language models that have already been fine-tuned and extended. To assess the performance of this step, we conduct empirical evaluation on our fine-tuned LLaMA2-256k model. This model is further extended to 512k, 1024k, and 2048k contexts using both PI and YaRN methods. As indicated by the results in Table~\ref{tbl:secondsearch}, our non-uniform positional interpolation method is able to maintain stable perplexity scores. Conversely, applying PI and YaRN results in a noticeable and rapid increase in perplexity as the extension factor grows.

\textbf{Improving recovery performance at reduced context lengths.} To address the degradation in performance observed at shorter context lengths, we reconfigure the RoPE factors for LongRoPE-2048k using our search algorithm. In particular, we restrict the upper limit of the scale factors considered during the search to promote reduced interpolation when evaluating 4k and 8k context sizes. Table~\ref{tbl:shortperformance} presents a comparison of perplexity for LongRoPE-LLaMA2-2048k on Proof-pile at 4k and 8k tokens, as well as the mean accuracy across LLM benchmarks. The findings reveal substantial gains in performance at these shorter context lengths.

\textbf{Analysis on the two forms of non-uniformities}. In this section, we investigate the individual contributions of the two types of non-uniformity by conducting ablation studies. Specifically, two experimental setups are considered: (i) we expand LLaMA2-7B to sequence lengths of 16k and 32k using multiple strategies—including PI, tuning only the RoPE dimension, and simultaneously optimizing both non-uniformities; (ii) we further extend a fine-tuned LLaMA2 model trained on 256k-length sequences to 2048k, employing the analogous procedures. All perplexity metrics are reported without further fine-tuning. As summarized in Table~\ref{tbl:searchdimension}, incorporating non-uniformity in the RoPE dimension yields a notable decrease in perplexity compared to the linear interpolation baseline used in PI. Adjusting for non-uniformity in token position leads to significant improvements at 16k and 32k, but this advantage diminishes at 2048k—potentially attributable to the exceptional length of the sequences. In cases where only the early token positions are retained without interpolation, the benefit disappears; future work will further explore this issue.

\section{Related Works}

In addition to methods based on position interpolation, this section discusses related works of other approaches. 

\textbf{Retrieval-based} methods rely on external memory modules to store distant contextual information, coupled with retrieval mechanisms that fetch pertinent documents during inference~\cite{longllama,wang2023augmenting,borgeaud2022improving}. Such systems often necessitate deliberate architectural changes to the LLMs. In comparison, our approach adopts a more streamlined methodology, requiring only minimal adjustments to positional embeddings. Furthermore, it supports a wider range of long context applications beyond document retrieval, including long document summarization and few-shot learning.

\textbf{Attention-based context window extensions}. In addition to methods involving positional embedding interpolation, several studies have increased the effective context window size within the limits of the base LLM’s window by innovatively adjusting the attention mechanism~\cite{han2023lminfinite, streamingllm,parallelwindow}. The central approach here is to address the exponential growth of attention computation at new positions by introducing specialized attention masks. These techniques are complementary to positional interpolation strategies.

\textbf{Fine-tuning based approaches} concentrate on optimizing the fine-tuning process of pre-trained LLMs equipped with adjusted position embeddings to support extended context windows. Approaches such as Code LLaMA~\cite{codellama}, LLaMA2 Long~\cite{xiong2023effective}, and ScaledRoPE~\cite{liu2023scaling} adopt a strategy of assigning an exceptionally large base value to RoPE and subsequently performing fine-tuning on sequences of the desired length. In contrast, our approach is adaptable to a wide spectrum of target sequence lengths and is capable of reaching lengths exceeding 2M. With the growing computational demands of fine-tuning on extensive contexts (e.g., exceeding 128k tokens), recent methods like LongLoRA~\cite{longlora} and PoSE~\cite{zhu2023pose} have emerged to address these resource challenges. Our technique is complementary to such efficient fine-tuning frameworks.

\section{Conclusion}

In this study, we introduce LongRoPE, a technique that significantly increases the context window of LLMs up to an extraordinary 2048k, all while preserving their effectiveness on tasks with the standard shorter context. Our approach leverages two types of non-uniform characteristics in RoPE positional embedding, employing an efficient evolutionary search to identify optimal configurations. This methodology yields dual advantages: it serves as a robust initialization point for subsequent fine-tuning and allows for an 8$\times$ extension of the context window even without additional fine-tuning. Furthermore, we develop a progressive extension strategy, utilizing LLMs fine-tuned with a 256k context window to systematically expand capacity to 2048k without the need for further fine-tuning. Comprehensive empirical results demonstrate the robustness and utility of LongRoPE. We anticipate that the LongRoPE-2048k models will open up possibilities for novel long-context use cases and catalyze future innovations in the field.

\section*{Broader Impacts} This research aims to contribute to progress within the discipline of Machine Learning. While our efforts may have various societal implications, we do not identify any particular effects that warrant explicit emphasis in this section.

	\balance

\end{document}